\section{Model Development}

After establishing the official MedMNIST ResNet-50 benchmark as an external
reference point, the first model trained in this work is a custom small CNN
trained from scratch which is used as a base model. Its architecture is shown
in Figure~\ref{fig:baseline224-smallcnn-architecture}. It consists of five
$3 \times 3$ convolutional layers. The first two layers map the RGB input from
three channels to 64 feature channels, followed by max pooling. The next two
convolutions increase the representation to 128 channels, again followed by max
pooling. A final convolution maps the representation to 256 channels. Every
convolution is followed by batch normalization and a SiLU activation. The
spatial feature map is then reduced by adaptive average pooling, followed by
dropout with probability $0.2$ and a linear classifier for the nine PathMNIST
classes. In total, the model has 558,921 trainable parameters.
\begin{figure}[htbp]
    \centering
    \begin{tikzpicture}[
        node distance=0.55cm,
        block/.style={
            draw,
            rounded corners,
            align=center,
            minimum width=2.3cm,
            minimum height=0.85cm,
            font=\small
        },
        arrow/.style={-{Latex[length=2mm]}, thick}
    ]
        \node[block] (input) {Input\\$224 \times 224 \times 3$};
        \node[block, right=of input] (conv1) {Conv block 1\\$3 \rightarrow 64 \rightarrow 64$};
        \node[block, right=of conv1] (pool1) {Max pool\\$112 \times 112$};
        \node[block, right=of pool1] (conv2) {Conv block 2\\$64 \rightarrow 128 \rightarrow 128$};
        \node[block, below=of conv2] (pool2) {Max pool\\$56 \times 56$};
        \node[block, left=of pool2] (conv3) {Conv block 3\\$128 \rightarrow 256$};
        \node[block, left=of conv3] (gap) {Adaptive\\average pool};
        \node[block, left=of gap] (head) {Dropout $0.2$\\Linear $256 \rightarrow 9$};

        \draw[arrow] (input) -- (conv1);
        \draw[arrow] (conv1) -- (pool1);
        \draw[arrow] (pool1) -- (conv2);
        \draw[arrow] (conv2) -- (pool2);
        \draw[arrow] (pool2) -- (conv3);
        \draw[arrow] (conv3) -- (gap);
        \draw[arrow] (gap) -- (head);
    \end{tikzpicture}
    \caption{Architecture of the standalone $224 \times 224$ baseline CNN.
    Each convolution is followed by batch normalization and a SiLU activation.}
    \label{fig:baseline224-smallcnn-architecture}
\end{figure}


We choose this smaller architecture instead of a ResNet because a ResNet
contains far more parameters and is more expensive to train on $224 \times 224$
inputs, while the small CNN can be trained and inspected more quickly. It is
also easier to analyze layer by layer, which is important because this work aims
to understand which visual structures are used for the cancer-related classes.
The model therefore serves as a transparent base model before adding more
specialized improvement strategies.

\subsection{Base Model Performance}%TODO: Why these hyperparameters?
Our base model is trained on the $224 \times 224$ PathMNIST dataset for 10
epochs with a batch size of 32. Optimization is performed with AdamW using a
learning rate of $0.001$ and weight decay of $0.0001$. A cosine learning-rate
schedule is used over the training run. To improve robustness, the training
configuration applies histology-specific data augmentation, PathMNIST
normalization, label smoothing of $0.03$, and mixup with $\alpha = 0.05$.

On the test set, the base model reaches an overall accuracy of $0.9057$, and an
overall $F_2$ score of $0.8850$. For the epithelium class, the model performs
strongly. It achieves a recall of $0.9886$ and an $F_2$ score of $0.9623$,
correctly classifying 1,219 of 1,233 test samples from this class. In contrast,
stroma also is a major weakness on our base model. Its recall is only $0.5202$
and its $F_2$ score is $0.5688$, with only 219 of 421 test samples classified
correctly.

From the confusion matrix shown in Figure~\ref{fig:baseline224-confusion-matrix}
we can infer the weak performance on stroma is not due to a uniform failure
across all classes. Instead, the base model classifies many true stroma samples
to certain classes, especially debris with 36 samples, smooth muscle with 78
samples, and epithelium with 77 samples.
\begin{figure}[htbp]
    \centering
    \includesvg[inkscape=false,inkscapelatex=false,width=0.6\textwidth]{figures/baseline224_confusion_matrix}
    \caption{Test-set confusion matrix of the standalone $224 \times 224$
    small CNN baseline. Rows correspond to ground-truth labels and columns
    correspond to predicted labels.}
    \label{fig:baseline224-confusion-matrix}
\end{figure}



\subsection{Feature Visualization and Interpretability}
To inspect why the base model performs well on epithelium but poorly on stroma,
we analyze what it has learned internally. For that, activation maps from
intermediate convolutional layers are visualized for both cancer-related
classes which can be seen in Figure~\ref{fig:baseline224-feature-maps}.
The feature maps indicate that early layers capture local texture and edge-like
structures, while deeper layers produce more spatially selective activations.
For the epithelium sample, several deeper feature maps appear to emphasize
glandular or boundary-like structures, whereas the stroma example shows more
diffuse textural activations. This suggests that the CNN learns
morphology-related representations, although feature-map inspection alone is
insufficient to establish that the learned features are reliable.
\begin{figure}[htbp]
    \centering
    \includesvg[inkscape=false,inkscapelatex=false,width=0.6\textwidth]{figures/baseline224_cancer_feature_maps}
    \caption{Layer-wise activation maps for a colorectal adenocarcinoma
    epithelium sample from the standalone $224 \times 224$ baseline CNN.}
    \label{fig:baseline224-cancer-feature-maps}
\end{figure}


% Grad-CAM highlights the image regions that contribute most strongly to the
% final class prediction.
To further assess whether the base model bases its predictions on meaningful
tissue regions, Grad-CAM visualizations were generated for representative
samples of the two cancer-related classes, which can be seen in
Figure~\ref{fig:baseline224-gradcam}.
For the stroma sample, the model activates several distributed regions across
the fibrous tissue structure. The highlighted areas roughly correspond to
regions containing elongated stromal texture and scattered nuclei. However, the
activation is relatively diffuse, which indicates that the decision may be based
more on general texture patterns than on sharply localized morphological
structures, therefore inducing a potential explanation for the weak $F_2$ score
of this class. For the epithelium sample, the Grad-CAM heatmap shows stronger
and more localized activations. Several high-activation regions are located
close to dense epithelial structures and tissue boundaries. Compared to the
stroma sample, the model appears to focus more clearly on spatially constrained
tissue regions, suggesting that the learned representation for epithelium may be
more meaningful.
\begin{figure}[htbp]
    \centering
    \includesvg[inkscape=false,inkscapelatex=false,width=\textwidth]{figures/baseline224_cancer_gradcam}
    \caption{Grad-CAM visualizations for cancer-associated stroma and colorectal
    adenocarcinoma epithelium samples.}
    \label{fig:baseline224-gradcam}
\end{figure}


\subsection{Expert Classifiers}
As can be seen in the confusion matrix in
Figure~\ref{fig:baseline224-confusion-matrix}, a substantial number of true
stroma samples are classified as debris, smooth muscle, or epithelium which is
the reason for the weak $F_2$ score. Together with the fact that the model's
Grad-CAM activations for stroma are relatively diffuse, we can hypothesize that
the baseline CNN is too conservative when predicting stroma.

This observation is further supported by the confidence distribution of true
stroma samples within their predicted classes. As shown in
Figure~\ref{fig:baseline224-stroma-confidence}, true stroma samples that are
incorrectly assigned to debris, smooth muscle, or epithelium often receive only
moderate prediction confidence. This indicates that at least part of the stroma
false negatives occur in uncertain decision regions rather than in cases where
the baseline model is highly confident. Therefore, we can use the confidence
distribution for triggering a second-stage correction step.
\begin{figure}[htbp]
    \centering
    \includesvg[inkscape=false,inkscapelatex=false,width=\textwidth]{figures/baseline224_stroma_confidence_distributions}
    \caption{Confidence distributions for true cancer-associated stroma samples
    under the base model.}
    \label{fig:baseline224-stroma-confidence}
\end{figure}


The overall confidence distributions for the competing non-stroma classes
provide an additional motivation for this strategy. As shown in
Figure~\ref{fig:baseline224-nonstroma-confidence}, many correctly predicted
non-stroma samples, especially for debris and epithelium, are assigned with very
high confidence. In contrast, the true stroma samples hidden inside these
predicted classes are more frequent in lower-confidence ranges. This separation
suggests that confidence can serve as a useful filtering criterion:
high-confidence predictions can be kept unchanged, while lower-confidence
predictions can be passed to a specialized binary classifier. The separation is
less clear for smooth muscle, where true stroma and non-stroma samples overlap
more strongly in confidence space. Therefore, the smooth-muscle expert may be
more difficult to apply without increasing the number of false-positive stroma
predictions.
\begin{figure}[htbp]
    \centering
    \includesvg[inkscape=false,inkscapelatex=false,width=\textwidth]{figures/baseline224_stroma_nonstroma_confidence_comparison}
    \caption{Confidence distributions for selected non-stroma predictions under
    the baseline $224 \times 224$ CNN. Blue bars show true cancer-associated
    stroma samples assigned to the predicted class, while red bars show samples
    from all other true classes assigned to the same predicted class.}
    \label{fig:baseline224-nonstroma-confidence}
\end{figure}


Based on these observations, a second-stage correction strategy is introduced.
For the main confusion pairs, binary expert classifiers are trained to
distinguish stroma from the respective competing class: debris versus stroma,
smooth muscle versus stroma, and epithelium versus stroma. These experts are
only consulted when the base model predicts one of the competing classes with
insufficient confidence lower than a threshold $\tau=0.9$. The resulting
architecture is shown in Figure~\ref{fig:cancer-expert}. Each expert uses the
same CNN architecture as the base model, but with a binary output layer, in
essence stroma vs. the competing class. Therefore, the expert classifiers are
only trained on samples from its own two ground-truth classes.
\begin{figure}[htbp]
\centering
\begin{tikzpicture}[
    scale=0.58, transform shape,
    >=Latex,
    font=\small,
    box/.style={
        draw,
        rounded corners,
        thick,
        align=center,
        minimum height=1.1cm,
        minimum width=2.8cm,
        fill=blue!5
    },
    decision/.style={
        draw,
        diamond,
        thick,
        align=center,
        aspect=2.2,
        inner sep=1.2pt,
        fill=orange!8,
        text width=4.2cm
    },
    output/.style={
        draw,
        rounded corners,
        thick,
        align=center,
        minimum height=1.0cm,
        minimum width=2.1cm,
        fill=green!5
    },
    outer/.style={
        draw,
        rounded corners,
        very thick,
        inner xsep=10pt,
        inner ysep=10pt
    }
]

% --------------------------------------------------
% Nodes
% --------------------------------------------------

% Input centered between upper and lower paths
\node (input) at (-3.8,0) [box] {Sample $x$};

% Upper path
\node (base) at (0,1.4) [box] {Base model $\mathcal M_B$};
\node (baseout) at (3.0,1.4) [output, minimum width=1.3cm] {$\hat y_B$};
\node (decision) at (7.9,1.4) [decision]
{$\hat y_B \in \mathrm{Experts},$\\[2pt]
$p_B(\hat y_B \mid x) < \tau$};

\node (finalbase) at (13.3,1.4) [output, minimum width=2.3cm]
{$\hat y = \hat y_B$};

% Lower path
\node (expert) at (7.9,-1.4) [box, minimum width=2.9cm]
{Expert $\mathcal M_E$};

\node (finalexpert) at (13.3,-1.4) [output, minimum width=2.3cm]
{$\hat y = \hat y_E$};

% Final output outside the system box
\node (final) at (16.7,0) [output, minimum width=1.4cm]
{$\hat y$};

% --------------------------------------------------
% Helper coordinates
% --------------------------------------------------

\coordinate (inbranch) at (-1.9,0);
\coordinate (merge) at (14.8,0);

% --------------------------------------------------
% Arrows
% --------------------------------------------------

% Sample branching to both paths
\draw[-, thick] (input.east) -- (inbranch);
\draw[->, thick] (inbranch) |- (base.west);
\draw[->, thick] (inbranch) |- (expert.west);

% Main upper path
\draw[->, thick] (base) -- (baseout);
\draw[->, thick] (baseout) -- (decision);

% Decision branches
\draw[->, thick] (decision) -- node[above] {no} (finalbase);
\draw[->, thick] (decision) -- node[right] {yes} (expert);

% Expert output
\draw[->, thick] (expert) -- (finalexpert);

% Merge both internal outputs into final output
\draw[thick] (finalbase.east) -| (merge);
\draw[thick] (finalexpert.east) -| (merge);
\draw[->, thick] (merge) -- (final.west);

% --------------------------------------------------
% Outer system box
% --------------------------------------------------

\node[outer, fit=(base)(baseout)(decision)(expert)(finalbase)(finalexpert)] (system) {};
\node[anchor=south west, font=\bfseries] at (system.north west)
{Cancer Expert};

\end{tikzpicture}

\caption{Architecture of the final cancer expert model.}
\label{fig:cancer-expert}
\end{figure}
